\section{Introduction}\label{sec:introduction}

Let $M\subseteq\mathcal B(H)$ be a von Neumann algebra on a complex Hilbert space, and let $A\subseteq M$ be a maximal abelian unital $*$-subalgebra, or \emph{masa}. A finite partition $P=(p_1,\ldots,p_r)$ in $A$ consists of mutually orthogonal projections with sum $1$. Its associated pinching is
\[
 S_P(x)=\sum_{i=1}^r p_i x p_i,\qquad x\in M.
\]
Paving asks how small this operator can be made after subtracting a diagonal element of $A$. Requiring a small norm for $S_P(x)-a$ itself is norm paving, which fails for general masas \cite[Theorem~3.3]{PV15}. The following strong-operator formulation instead permits a compression whose complement is arbitrarily small in the strong operator topology.

\begin{definition}[Strong-operator paving]\label{def:so-paving}
Let $x=x^*\in M$, $\eps>0$, and $r\ge1$ be an integer. We say that $x$ is \emph{$(\eps,r)$ so-pavable over $A$} if for every finite $F\subset H$ and every $\delta>0$ there are a partition $P=(p_1,\ldots,p_r)$ in $A$, an element $a=a^*\in A$ with $\|a\|\le\|x\|$, and a projection $q\in M$ such that
\[
 \|(1-q)\xi\|<\delta\quad(\xi\in F),
 \qquad \|q(S_P(x)-a)q\|\le\eps\|x\|.
\]
The integer $r$ is fixed before $F$ and $\delta$ are chosen.
\end{definition}

\begin{theorem}\label{thm:main}
For every $0<\eps<1$, set
\[
 r_\eps=\left\lceil400000000\eps^{-2}\right\rceil
              +\left\lceil4/\eps\right\rceil+1
          \le500000000\eps^{-2}.
\]
For every masa $A\subseteq M\subseteq\mathcal B(H)$ and every $x=x^*\in M$, the element $x$ is $(\eps,r_\eps)$ so-pavable over $A$. In particular, the number of projections is fixed before the inclusion, operator, and strong neighborhood. No separability or conditional-expectation assumption is imposed on the inclusion.
\end{theorem}

Theorem~\ref{thm:main} proves the quadratic so-paving conjecture of Popa and Vaes \cite[Conjecture~2.8(2)]{PV15}. Their definition of so-paving uses arbitrary strong neighborhoods of zero; Definition~\ref{def:so-paving} gives the equivalent finite-vector formulation \cite[Definition~2.2]{PV15}. For self-adjoint $x$, allowing a general diagonal element in their definition gives the same condition, since taking its real part preserves the norm bound and the compression estimate.

For countably decomposable expected masas, the uniform bound also yields norm paving in the associated ultrapower.

\begin{corollary}[Quadratic norm paving in Ocneanu ultrapowers]\label{cor:ultrapower-paving}
Let $A\subseteq M$ be a masa. Suppose that $A$ is countably decomposable and is the range of a normal conditional expectation from $M$. Let $\omega$ be a free ultrafilter on $\mathbb N$, and let $A^\omega\subseteq M^\omega$ be the associated inclusion in the Ocneanu von Neumann ultrapower. For every $0<\eps<1$ and every $x=x^*\in M^\omega$, there are a partition $p_1,\ldots,p_{r_\eps}\in A^\omega$ and an element $a=a^*\in A^\omega$ such that
\[
 \|a\|\le\|x\|,
 \qquad
 \left\|\sum_{i=1}^{r_\eps}p_i x p_i-a\right\|\le\eps\|x\|.
\]
Here $r_\eps$ is the same integer as in Theorem~\ref{thm:main}.
\end{corollary}

\begin{proof}
Theorem~\ref{thm:main} makes every self-adjoint element of $M$ $(\eps,r_\eps)$ so-pavable. The forward transfer in \cite[Theorem~2.7(3)]{PV15} gives the stated norm paving in $M^\omega$ with the same error and number of projections. Taking the real part of the diagonal element preserves both estimates.
\end{proof}

\subsection{History and significance}

The classical Kadison--Singer paving problem concerns the diagonal masa in $\mathcal B(\ell^2\mathbb N)$. Marcus, Spielman, and Srivastava resolved it by their mixed-characteristic-polynomial method \cite{MSS15}. For general masas, norm paving is much more restrictive: when $M$ has separable predual, it holds exactly when $M$ is type~I and $A$ is the range of a normal conditional expectation \cite[Theorem~3.3]{PV15}. Strong-operator paving retains a uniform partition count while allowing a compression that approaches the identity strongly.

Popa's $\mathrm{II}_1$-factor approach established paving for singular-masa ultraproducts \cite{Popa14}. Popa and Vaes then obtained so-paving bounds of order $\eps^{-4}$ for arbitrary masas in type~I algebras with separable predual, for Cartan masas in amenable algebras, and for the Cartan inclusions associated with essentially free ergodic probability-measure-preserving profinite actions of countable groups \cite[Proposition~3.2 and Theorem~4.1]{PV15}. In contrast, their theorem for tracial ultraproducts of singular masas in finite algebras gives a common norm-paving partition of order $\eps^{-2}$ for each finite family of centered self-adjoint contractions \cite[Theorem~5.1]{PV15}. Their later work gives the optimal uniform error for singular-masa ultraproducts \cite[Theorem~4.1]{PV16}. They explicitly identified the need to unite the Cartan and singular arguments \cite[Remark~2.9(ii)]{PV15}. On the finite-dimensional side, Ravichandran and Srivastava obtained asymptotically optimal multipaving bounds from mixed determinantal polynomials \cite{RS17}; these supply the finite input used here.

A previous combination result illustrates the quantitative issue. For an intermediate inclusion $A\subseteq N\subseteq M$ with a faithful normal tracial state, Popa and Vaes assume that $L^2(M\ominus N)$ has no nonzero $A$--$N$ subbimodule finitely generated on the right. Their bound uses at most $400\eps^{-2}$ times the so-paving count for $A\subseteq N$ at error $\eps/2$ \cite[Remark~5.3]{PV15}. Inserting their quartic Cartan bound gives order $\eps^{-6}$. The present argument treats the normalizer algebra through a nonsingular measured equivalence relation and uses a single partition for both terms in the normalizer decomposition of one operator. This retains quadratic order. The finite-family obstruction in \cite[Remark~5.2]{PV15} concerns the stronger requirement of a common partition for arbitrarily many separately prescribed operators in a finite ambient algebra.

The exponent two is forced already by the $L^2$ obstruction in a $\mathrm{II}_1$ factor from \cite[Proposition~5.4]{PV15}. Fix an integer $1\leq n<\eps^{-2}$, choose $\eps<\eta<n^{-1/2}$, and use the centered self-adjoint unitary $x$ from that proposition at threshold $\eta$. Write $\tau$ for the normalized trace and $\|b\|_2=\tau(b^*b)^{1/2}$. For every partition $P$ with $n$ pieces and every $a\in A$, trace orthogonality gives
\[
 \|S_P(x)-a\|_2^2=\|S_P(x)\|_2^2+\|a\|_2^2.
\]
If $\|a\|\leq1$, set $z=S_P(x)-a$, so $\|z\|\leq2$. Whenever $\|qzq\|\leq\eps$, the two discarded corner terms give
\[
 \|S_P(x)\|_2\leq\|z\|_2
 \leq\eps+4\tau(1-q)^{1/2}.
\]
In the trace representation, a strong neighborhood can force $\tau(1-q)<((\eta-\eps)/4)^2$. The displayed estimate would then contradict the proposition's lower bound $\|S_P(x)\|_2>\eta$, which holds for every such $P$. Thus its obstruction persists with the varying diagonals and compressions allowed here. The exponent in Theorem~\ref{thm:main} is optimal, although its numerical constant is not optimized.

Theorem~\ref{thm:main} controls a compression of the pinching of the original $x$. A separate expectation-free approximation-paving theorem instead supplies self-adjoint approximants $y$ near $x$ strongly, with $\|y\|\leq3\|x\|$, and norm-paves them with error measured against $\|y\|$ \cite[Theorem~1.1]{ApproxCompanion26}. Its arbitrary-masa conclusion has a finite $\eps$-dependent count but states no quadratic rate. The two proofs are independent.

\subsection{Outline of the proof}

Let $p\in A$ be the supremum of the supports of all normal positive functionals $\rho$ satisfying $\rho(bz)=\rho(zb)$ for $b\in A$ and $z\in M$. After the finite vector tests are fixed, we choose one such support $e\leq p$ so that $p-e$ is small on those vectors. When $e\ne0$, the corner $eMe$ carries a faithful normal state central over $eAe$. On $(1-p)M(1-p)$, a symmetric Haagerup $L^2$ estimate and two randomizations give a paving with $O(\eps^{-1})$ colors; a positive-cone overlap estimate converts the symmetric estimate into a compression of arbitrarily large prescribed state mass. The corner $p-e$ is discarded. Since $p$ and $e$ need not be central in $M$, the construction uses disjoint palettes on these corners so that pinching removes cross-corner terms and the counts add. A separable reduction on the faithful-state corner then permits measured-space methods. These reductions are proved in Section~\ref{sec:reductions}.

In the presence of this central state, let $N$ be the algebra generated by the partial normalizers of $A$. There are $\varphi$-preserving conditional expectations $E_A:M\to A$ and $E_N:M\to N$, and $A\subseteq N$ is Cartan. We use the decomposition
\[
 x-E_Ax=(E_Nx-E_Ax)+(x-E_Nx).
\]
Section~\ref{sec:kernels} proves a structural fact about the second term: for $y\in\ker E_N$, the positive map $f\mapsto E_A(y^*fy)$ on $A$ is represented by an atomless kernel. This is the input needed to remove coincidences in the later moment transfer.

The Cartan term is treated on a relative Bernoulli extension of the measured equivalence relation representing $A\subseteq N$. Each relation class receives independent site labels. We construct a measurable coloring there that has uniform conditional color marginals and makes the pinched matrix small on any prescribed finite-radius ball with arbitrarily high probability. This is Theorem~\ref{thm:bernoulli-coloring}. Its proof occupies Sections~\ref{sec:polynomials}--\ref{sec:coloring}.

The finite-dimensional input to this construction is the mixed determinantal bound of \cite[Theorem~25]{RS17}. For a partial coloring, a real-rooted polynomial uses both signs of the matrix and shifts determined by the assigned sites; the sum of its positive roots is the total cost. At the chosen shift, the upper mixed bound makes the all-unassigned cost zero, while the lower comparison for the two signs makes excessive norm on a fully assigned finite set force positive cost. Rooted spectral measures allocate this cost to vertices, and a single-site assignment inequality controls its increase. For a nonsingular relation, an invariant lift with an extra real coordinate makes mass transport available. A deletion estimate controls large weights uniformly, so joint cutoff limits give homogeneous costs on the original probability space. Random assignment times produce a partial coloring of arbitrarily small total cost for which any prescribed ball is fully assigned with high probability. Completing the remaining sites preserves the estimate on those balls. The uniform cutoff estimates and the order of conditioning on the assignment times are part of the construction.

To combine the two operator terms, we place $M$ and the Bernoulli extension of $N$ in their amalgamated free product over $N$. Uniform conditional color marginals give a free compression bound on $\ker E_N$, while the same color projections provide the Cartan estimate. Section~\ref{sec:transfer} transfers each fixed mixed scalar moment from this model to actual projections of $A$. Its proof uses finite-symbol evaluation, modular analytic approximation, and a shortest-pair argument in cyclic words. The atomless kernels make the remaining collision terms vanish.

Finally, Section~\ref{sec:completion} converts the local Cartan estimates to bounded approximants by clipping, controls continuous calculus for the resulting varying sequences, and transfers a sufficiently high fixed moment. A spectral cut produces the required projection $q$. All approximation parameters are chosen after the number of colors, and the corner assembly completes Theorem~\ref{thm:main}.

\subsection{Conventions}

We allow zero projections in a partition and use $1$ also for the identity of a corner when that corner is regarded as an algebra. Except when restoring scale in the final proof, the operators to be paved are self-adjoint contractions. Maximal abelianness implies that $A$ is weakly closed and that $A'\cap M=A$: its weak closure is abelian, and any self-adjoint element of the relative commutant could otherwise be adjoined. In particular, corners by projections in $A$ are again masas in the corresponding corners of $M$.

For a normal positive functional $\varphi$, write $\|z\|_\varphi=\varphi(z^*z)^{1/2}$. When $\varphi$ is a faithful normal state, its centralizer is the algebra of elements fixed by its modular automorphism group. Measurable statements are understood modulo null sets. Whenever a countable nonsingular relation is present, a null exceptional set can be enlarged to its null saturation.
